\subsection{The local extension and its packets}

\paragraph{Statement C1 (Proved; correction to the centralizer reading).}\label{res:C1}
Let $U=X\setminus S$ be a hyperbolic curve over a number field $K$, let $x\in S(K)$, and fix a lift $c$ of $x$ to the completed geometric universal cover. With the chosen base point, the local arithmetic group is
\[
 D_c=\Norm_{E_U}(I_c),\qquad
 1\longrightarrow I_c\simeq\Zhat(1)\longrightarrow D_c
   \longrightarrow\Gamma_K\longrightarrow1.
\]
Its conjugation action on $I_c$ is $\chi_K$. In particular
\[
 \Cent_{E_U}(I_c)=D_c\cap p^{-1}(\Gamma_K^\circ).
\]
A local parameter supplies a splitting. The set of $I_c$-conjugacy classes of splittings is a torsor under
\[
 H^1_{\mathrm{cont}}(\Gamma_K,\Zhat(1))
   \simeq\varprojlim_{n\mid m} K^*/K^{*n}.
\]
These identifications describe the torsor; identifying it with its acting group additionally chooses a splitting.
\passage{Section 3, manuscript pp.~26--27, scans 53--55; formula (2) and the first formula (3).}
\begin{proof}
We use the following exact source input. Stix~\cite[Proposition 11 and Corollary 12]{Stix2008}, for a hyperbolic curve over a characteristic-zero field, proves that distinct lifted cusps have inertia groups with trivial intersection and that the geometric and arithmetic normalizers of $I_c$ are $I_c$ and $D_c$, respectively. The identification of local inertia with $\Zhat(1)$ is the characteristic-zero tame description: over an algebraically closed residue field the covers of a punctured strictly henselian smooth disc are the compatible covers $t=u^n$. Conjugating $u\mapsto\zeta_nu$ by a coefficient automorphism sends $\zeta_n$ to $\zeta_n^{\chi_n(\sigma)}$. Passing to the inverse limit gives the displayed action and the formula for the centralizer. A centralizer element lies in the normalizer; within the normalizer its action is trivial exactly when every $\chi_n$ equals $1$.

Choose the compatible roots $t^{1/n}$ and let $\Gamma_K$ act on coefficients while fixing those roots. This constructs a splitting of the local extension; restriction to the finite extensions in the chosen algebraic closure gives the indicated embedding into $E_U$. Any second splitting has the unique form $s_a(\sigma)=a(\sigma)s_0(\sigma)$, with $a$ a continuous $1$-cocycle in $I_c$. Conjugation by an element of $I_c$ changes $a$ by a coboundary. This proves the torsor assertion without choosing an origin for it.

For each $n$, the exact sequence $1\to\mu_n\to\overline K^*\xrightarrow{n}\overline K^*\to1$ gives $H^1(K,\mu_n)=K^*/K^{*n}$. Here the vanishing used for $H^1(K,\overline K^*)$ has the following elementary proof. A continuous cocycle descends to a finite Galois extension $L/K$. Linear independence of the automorphisms of $L$ supplies $b\in L$ for which $\sum_\sigma a_\sigma\sigma(b)\ne0$; applying an automorphism to this sum expresses the cocycle as a coboundary. At the inverse-limit step, the finite groups $\mu_n(K)$ have the Mittag--Leffler property, so the possible $\varprojlim^1 H^0$ term is zero. Equivalently, the finite nonempty sets of compatible representatives of a fixed compatible family of cohomology classes have nonempty inverse limit. Thus continuous $H^1$ with coefficients in $\Zhat(1)$ is the asserted limit of Kummer groups.

For a number field, $\chi_K(\Gamma_K)$ is open in $\Zhat^\times$: $K\cap\Q(\mu_\infty)$ is finite over $\Q$, and restriction identifies this image with the subgroup fixing that intersection. It is nontrivial. Consequently the literal centralizer in scan 54 cannot be the middle group of its extension onto all of $\Gamma_K$. We retain the word visible in the scan and record the required mathematical replacement in the readings file.
\end{proof}

\paragraph{Statement C2 (Proved; the packets are embedded and disjoint).}\label{res:C2}
For hyperbolic $U/K$, the map from $I_c$-conjugacy classes of splittings of $D_c\to\Gamma_K$ to $\pi_U$-conjugacy classes of sections of $E_U\to\Gamma_K$ is injective. Its image is disjoint from the point sections of $U(K)$ and from the packet of any other rational cusp. These assertions include type $(0,3)$ and remain true after each finite extension of $K$.
\passage{Section 3, manuscript pp.~26--27, scans 54--55, including the exception left for type $(0,3)$.}
\begin{proof}
First use the proved centralizer theorem \Rtwo. A local splitting $s$ has $C(s)=I_c$. If two local splittings satisfy $s'=hsh^{-1}$ with $h\in\pi_U$, then $I_c=C(s')=hC(s)h^{-1}=hI_ch^{-1}$. The geometric normalizer formula in C1 puts $h$ in $I_c$. This proves injectivity. For splittings at different cusps, the same equality identifies their lifted inertia groups after conjugation, hence their images in $S$; the two cusps must coincide. A point section has zero cyclotomic centralizer by \Rtwo\ and the separation argument below, whereas a local splitting has the nonzero group $I_c$.

Here is a second proof that also checks the separation used in the preceding sentence. Stix~\cite[Theorem 14(1)--(2)]{Stix2008} assumes injectivity of the point-to-section map for proper curves of genus at least two over every finite extension of the base field. Under exactly that hypothesis it proves all the packet assertions in C2. Its proof first replaces $U$ by a finite étale neighbourhood of the section whose compactification has genus at least two. Two sections coming from different points of the compactification are then separated there. Repeating over all finite neighbourhoods separates different lifted cusps over the same cusp. An element conjugating two local splittings therefore stabilizes the selected lifted cusp, and lies in $I_c$.

For completeness the arithmetic hypothesis follows as follows. Given distinct rational points $a,b$ of a smooth proper curve $Y$ of positive genus, the Abel--Jacobi map based at $a$ sends $b$ to $[b-a]\in J_Y(K)$. This class is nonzero: a function with divisor $b-a$ would have degree one as a map to $\mathbf P^1$, forcing $Y$ to have genus zero. Equality of the two sections would make $[b-a]$ divisible by every positive integer: pull back the multiplication-by-$n$ covering of $J_Y$, and read the equality of the two fibre torsors. The no-divisible-elements statement \Arith\ excludes this. The argument works over every finite extension. The higher-genus neighbourhood can be constructed explicitly over the algebraic closure and then descended by a characteristic refinement and the section. In genus zero, carry three punctures to $0,1,\infty$ and pull back the Fermat quartic; its compactification has genus three. In genus one, first take an etale double cover of the compactification, so that its open part has at least two cusps $P,Q$. Choose a degree-one line bundle $L$ with $L^{\otimes2}\simeq\mathcal O(P+Q)$; it exists because multiplication by two on the elliptic curve's degree-zero Picard group is surjective. The double cover defined by the section with divisor $P+Q$ is branched at these two cusps and its compactification has genus two by Riemann--Hurwitz. In compact genus at least two the identity cover suffices. Each constructed cover is etale on the relevant open curve. Take its Galois closure and an open characteristic subgroup contained in its subgroup. The new compactification dominates the constructed one, so Riemann--Hurwitz preserves genus at least two. The subgroup together with the section descends this cover.

The published Esnault--Hai statement~\cite[Theorem 5.5]{EsnaultHai2008} concerns a smooth projective $X$ over a finitely generated extension of $\Q$, a nonempty open $U$, and a neutral fibre functor for finite connections on $U$. For $g(X)\ge1$, or $g(X)=0$ with $U$ hyperbolic, it bounds the number of rational points in the corresponding normalized universal cover by one. Its Lemma 5.6 supplies a positive-genus finite level in genus zero over any characteristic-zero field. The proof just given is the group-theoretic version of this positive-genus reduction. The attached older preprint has a different item 5.6; the published lemma has been opened and checked.
\end{proof}

\paragraph{Statement C3 (Refuted as written in type $(0,2)$; exact replacement).}\label{res:C3}
For $U=\Gm/K$, both cuspidal packets are the whole set
$H^1(K,\Zhat(1))$ of section classes. In particular they contain every point section. For an affine curve with nontrivial geometric fundamental group, the disjointness conclusion of scan 54 holds for every hyperbolic type and fails for type $(0,2)$. In types $(0,0)$ and $(0,1)$ the geometric group is trivial. For a proper genus-one curve admitting a section, its cyclotomic centralizer is zero over a number field.
\begin{proof}
The covers $u^n=t$ of $\Gm$ identify $\pi_U$ with $\Zhat(1)$. The point $1$ splits the arithmetic extension, giving $E_U=\Zhat(1)\rtimes\Gamma_K$. The local group at $0$ is this entire group, and the coordinate $t^{-1}$ at infinity gives the same group with inertia generator inverted. Thus each map from local splittings to global splittings is surjective, indeed bijective on conjugacy classes. A point $a\in K^*$ gives its Kummer class in precisely the same group. Since the geometric kernel is abelian, the conjugation action of every splitting on it is $\chi_K$ and $C(s)=\Zhat(1)$ for every splitting.

A sphere and a once-punctured sphere are simply connected over $\mathbf C$; finite-cover comparison identifies their geometric étale groups with $1$. A torus has fundamental group $\Z^2$, whose profinite completion is the Tate module of its Jacobian. In the genus-one case changing a section conjugates by this abelian kernel and does not change its action; \Arith\ gives zero invariants. The finite-cover comparison used here is SGA~1, Exposé XII, Theorem 5.1: analytification induces an equivalence between finite étale covers of a complex scheme locally of finite type and finite étale analytic covers. The comparison with the complex fibre is valid in characteristic zero, as in the same exposé and the specialization statement of Exposé XIII; all our curves descend to finitely generated subfields embedded in $\mathbf C$.
\end{proof}

\paragraph{Statement C4 (Proved; the precisions without section surjectivity).}\label{res:C4}
Over a number field, the section of any rational point of a hyperbolic curve has $C(s)=1$. Every cuspidal section has $C(s)=I_c$ for exactly one lifted cusp $c$. The same formulas hold for germs over finite extensions and for restriction to an arbitrary open subgroup of the appropriate cyclotomic kernel. These formulas alone impose no existence statement for sections whose centralizer is trivial.
\passage{Section 3, manuscript p.~27, scans 55--56; the second formula (3) and the marginal precisions to Conjecture A.}
\begin{proof}
The local formula and its uniqueness are \Rtwo; C2 excludes a point section from every local packet. The open-subgroup statement is the germ version in Part A. More explicitly, an open subgroup of $\Gamma_K^\circ$ is $\Gamma_{F'}$ for a finite extension $F'/K(\mu_\infty)$, and the arithmetic vanishing and packet separation continue over that field. The normality of $\Gamma_K^\circ$ then makes its uniquely detected cusp invariant under the whole section. None of these deductions asserts that a section with zero centralizer is in the image of $U(K)$.
\end{proof}

\subsection{Dévissage and the quantifier over all open subgroups}

\paragraph{Statement C5 (Proved; exact sequences for elementary fibrations).}\label{res:C5}
Let $Y\to B$ be an elementary fibration in hyperbolic curves in characteristic zero, with $B$ obtained by a finite tower of such fibrations from a point. The elementary hypothesis means a smooth proper curve compactification whose omitted divisor is finite étale over the base. At compatible geometric base points, the maps induced by the fibre inclusion and by the fibration give an exact sequence of profinite groups
\[
 1\longrightarrow\pi_1(Y_{\bar b})\longrightarrow\pi_1(Y_{\overline K})
 \longrightarrow\pi_1(B_{\overline K})\longrightarrow1.
\]
For $b\in B(K)$ there is also a natural isomorphism of extensions
\[
 E_{Y_b}\ \xrightarrow{\ \sim\ }\ E_Y\times_{E_B}s_b(\Gamma_K).
\]
\begin{proof}
We give the part of the proof which prevents a loss of injectivity on completion. Bresciani, Lemma A.22, supplies invariance of geometric fundamental groups under extension of algebraically closed characteristic-zero fields. After descending to a field of definition and embedding it in $\mathbf C$, the smooth proper family and its finite collection of disjoint marked sections are locally trivial as differentiable pairs. To see this locally, lift tangent vector fields from a relatively compact coordinate ball in the base to the total space, choosing the lifts tangent to the marked divisor; a partition of unity patches the lifts, and properness supplies the flows throughout a smaller ball. The flows trivialize the pair, and hence its open complement. Hyperbolic complex curves are aspherical. The homotopy lifting property and the long exact homotopy sequence therefore prove inductively that the base and total space are aspherical and that
\[
 1\to P\to G\to B_0\to1
\]
is exact for their discrete topological fundamental groups.

The goodness input can be stated precisely: Bresciani~\cite[Lemmas 10.9 and 10.11]{Bresciani2021} proves goodness and residual finiteness for elementary anabelian topological stacks, with goodness preserved under extensions whose kernel has finite cohomology with finite coefficients and is finitely generated. For open curve fibres the kernel is a finitely generated free group. Its discrete and completed free-group resolutions have length one; in degree one the comparison identifies derivations by their values on the same free generators, and in degrees at least two both cohomologies vanish. This proves goodness and the finite-cohomology property in this case. Lemma 10.10 of the same source proves that finite-index subgroups of a good group remain good, as needed below. Thus the identical induction includes the open elementary fibrations used here. We now justify the resulting profinite exactness rather than treating it as a formal property of completion.

Take a characteristic finite-index subgroup $N$ of $P$ and put $A=P/N$. The quotient $G/N$ has finite normal subgroup $A$. Its subgroup centralizing $A$ has finite index, meets $A$ in $Z(A)$, and is a central extension of a finite-index subgroup $B_1\subset B_0$ by the finite abelian group $Z(A)$. Goodness of $B_1$ identifies its $H^2$ with continuous $H^2$ of $\widehat B_1$. A continuous cocycle with finite coefficients factors through a finite quotient. The central-extension class thus restricts to zero on a finite-index subgroup $B_2$, and over $B_2$ a splitting supplies a subgroup disjoint from $Z(A)$. Its preimage in $G$ has finite index and intersects $P$ inside $N$. Characteristic finite-index subgroups are cofinal in $P$: intersect the finitely many subgroups of a fixed index bound in a finitely generated group. Hence the full profinite topology on $P$ is induced from $G$, proving $\widehat P\hookrightarrow\widehat G$. The quotient by the closure of $P$ is $\widehat B_0$ because its finite quotients are exactly the finite quotients of $G/P$. This proves the short exact sequence after completion. SGA~1 XII.5.1 identifies these completions with the geometric étale groups.

For the second assertion the morphism of schemes $Y_b\to Y$ induces the displayed map, with identity on the quotient $\Gamma_K$. Its restriction to geometric kernels is the identification of the kernel just proved, so injectivity and surjectivity follow by lifting an element to the quotient and then correcting it in the kernel. All maps are continuous; a continuous bijection of compact profinite groups is a homeomorphism.
\end{proof}

% Optional supporting paragraph for C5; root chooses its insertion point.
An independent exact reference, with no goodness or embedding assumption,
is Hoshi~\cite[Proposition 2.4(i),(ii), pp.~169--172]{B_Hoshi2014}.
His base $S$ is connected, noetherian, separated and normal over a field
of characteristic zero, and his relative hyperbolic curves have smooth
proper compactifications and finite \emph{\'etale} cusp divisors.
Part (i), with $n=1$, $m=0$, $S=B_{\overline K}$ and
$X=Y_{\overline K}$, gives precisely the displayed geometric exact
sequence, including injectivity on the left. Part (ii), with
$S=B$, $X=Y$ and $T=\operatorname{Spec}K\xrightarrow{b}B$, gives
$E_{Y_b}\simeq E_Y\times_{E_B}\Gamma_K$, with the two maps supplied
by the fibre inclusion and the structure map. Compatible geometric
base paths turn his outer isomorphisms into the representatives used
here.

\paragraph{Statement C6 (Proved; point sections of elementary anabelian varieties).}\label{res:C6}
Let $U/K$ be an elementary anabelian variety over a number field and $u\in U(K)$. For every open $\Omega\subset\Gamma_K^\circ$,
\[
 \Cent_{\pi_U}(s_u(\Omega))=1.
\]
Moreover every morphism $f:V\to U$ of nonempty connected finite-type schemes over $\Q$ induces a representative $f_*:E_V\to E_U$ such that
\[
 \Cent_{E_U}(f_*(H))=\Cent_{\pi_U}(f_*(H))=1
 \quad\text{for every open }H\subset E_V^\circ.
\]
Changing geometric base paths conjugates the maps and their centralizers and does not affect these equalities.
\passage{Section 2, manuscript pp.~23--24, scans 48--49, condition b) and its asserted reduction to a point.}
\begin{proof}
Induct along the fibration tower. An element centralizing $s_u(\Omega)$ maps to an element centralizing the point section at the image $b$ in the base. By induction its image is $1$. C5 places it in the geometric group of the fibre over $b$, and identifies the action on that kernel with the point section of the fibre at $u$. C4 for that hyperbolic fibre makes the element $1$. The zero-dimensional base case has trivial geometric group. All comparison maps are the fibre inclusions and fibration projections of C5, so the induction uses injectivity precisely where needed.

Now let $H$ be open in the closed subgroup $E_V^\circ$. There is an open normal subgroup $N\triangleleft E_V$ with $N\cap E_V^\circ\subset H$: an identity neighbourhood in the subspace topology contains the intersection with an open normal neighbourhood in $E_V$. Let $V_N\to V$ be the corresponding connected finite étale cover. It has a closed point $v$ whose residue field $k$ is a number field, since it is nonempty of finite type over $\Q$. Choose base paths so that the induced point homomorphism $s_v:\Gamma_k\to E_V$ has image in $N$. Its cyclotomic kernel has image in $N\cap E_V^\circ\subset H$. Under $f_*$ this is the restriction of the point section at $f(v)$, with arithmetic domain $\Gamma_k$. The first part therefore gives $\Cent_{\pi_U}(f_*(H))=1$.

An element of the arithmetic centralizer projects to an element of $\Gamma_\Q$ centralizing the image of $\Gamma_k^\circ$. Part B proves that centralizer trivial. Thus the arithmetic element lies in $\pi_U$, where it is already zero. This proves the other equality. The argument displays a point after choosing $H$, so it preserves the manuscript's universal quantifier over all open $H$. Regularity of $V$ is unnecessary for this direction.
\end{proof}

\paragraph{Statement C7 (Proved equivalence of predicates; surjectivity remains unresolved over number fields).}\label{res:C7}
For a hyperbolic curve $U/K$, let $\mathcal S(U/K')$ denote section classes over a finite extension $K'/K$, and $\mathcal S_{\mathrm{nc}}(U/K')$ its noncuspidal part. Condition b) of scan 48 for $V=\operatorname{Spec}K'$ holds for exactly $\mathcal S_{\mathrm{nc}}(U/K')$. Consequently the unresolved map
\[
 U(K')\ \lhook\joinrel\longrightarrow\ \mathcal S_{\mathrm{nc}}(U/K')
\]
has the same surjectivity defect as both the point case of the fundamental conjecture and Conjecture A. This is an equality of the two requested assertions, without assuming either one.
\begin{proof}
C4 and \Rtwo\ identify trivial centralizer on every open subgroup of $\Gamma_{K'}^\circ$ with noncuspidality. C2 puts all point classes in this subset, and the Abel--Jacobi and finite-cover argument proves the displayed injectivity. Every other class is cuspidal exactly when the displayed map is onto the noncuspidal subset. The same statement, read as a claim that every class satisfying b) arises from a point, is the point case of scan 48. These implications apply separately over every finite $K'$. The smallest instance left by this argument is a noncuspidal section of $\mathbf P^1\setminus\{0,1,\infty\}$ over a number field that has not been identified with a point. The calculations in Part D separate a specific pair of point sections; they do not provide a point for an arbitrary class in this unresolved set.
\end{proof}

% Root splice: replaces Statement C8, not the surrounding subsections.
% Root splice: replaces Statement C8, not the surrounding subsections.
\paragraph{Statement C8 (Proved: boundary nonvanishing in every dimension;
local injectivity refuted).}\label{res:C8}
Let $U/K$ be an elementary anabelian variety, namely a finite tower of
relative hyperbolic curves over a field of characteristic zero, and let
$U\hookrightarrow X$ be a smooth proper compactification with
normal-crossings boundary. For a rational boundary point $x$, write
$I(x)$ for the nonempty set of geometric local branches at $x$, and put
$T=\widehat{\mathbf Z}(1)$. The geometric group of the punctured strict
henselian neighbourhood is $T^{I(x)}$, and its natural map
\[
 j_x:T^{I(x)}\longrightarrow\pi_1(U_{\overline K})
\]
has nonzero image. More precisely, its restriction to the diagonal
$d_x:T\to T^{I(x)}$ has nonzero image. Every section through the
punctured henselian neighbourhood therefore has nontrivial cyclotomic
centralizer. When all local branches are rational, that centralizer
contains the whole image of $j_x$; in general it contains the nonzero
image of $j_xd_x$. These assertions hold over every finite base extension
and hence for germs. The assertion that $j_x$ is injective for every
such compactification is false, already on a blown-up product surface.
\passage{Section 4, manuscript pp.~63--65, scans 128--131, especially
the withdrawn assertion $1^\circ$ and the margin at scan 130.}
\begin{proof}
The punctured henselian neighbourhood is
$U_{(x)}=U\times_X\operatorname{Spec}\mathcal O^h_{X,x}$.
Its local exact sequence and action are computed in
Carlson--Corwin~\cite[Proposition 3.4]{B_CarlsonCorwin2026}:
the geometric kernel is
\[
 \operatorname{Hom}\bigl(\mathbf Z[I(x)],T\bigr)=T^{I(x)},
\]
with Galois action obtained by permuting the branches and by the
cyclotomic action on $T$. The natural morphism $U_{(x)}\to U$ supplies
$j_x$ and makes it equivariant for every chosen local section. If the
branches are rational, the cyclotomic kernel fixes this entire local
geometric group. Without that restriction it still fixes the diagonal,
because every permutation fixes a constant tuple. It remains to prove
that the diagonal survives the displayed map.

We use the following precise valuative input. Hoshi
\cite[Lemma 2.8, pp.~175--176]{B_Hoshi2014} treats a relative hyperbolic
curve $Y\to S$, where $S$ is connected, noetherian, separated and
normal over a characteristic-zero field. For a strictly henselian
discrete valuation ring $R$ over $S$, with fraction field $F$, a given
$S$-map $\operatorname{Spec}F\to Y$ extends over
$\operatorname{Spec}R$ exactly when its homomorphism
$\pi_1(\operatorname{Spec}F)\to\pi_1(Y)$ has trivial image.

This gives the same extension criterion for a hyperbolic polycurve over
an algebraically closed characteristic-zero field. Here is the
induction, including the maps involved. For a tower
$U=U_n\to U_{n-1}\to\cdots\to U_0=\operatorname{Spec}\overline K$
and $a:\operatorname{Spec}F\to U$ with trivial fundamental-group
image, its composite to $U_{n-1}$ also has trivial image. By induction
that composite extends to $b:\operatorname{Spec}R\to U_{n-1}$.
Regard $R$ as a ring over the normal base $U_{n-1}$ by $b$.
Hoshi's lemma applied to the last relative curve and to the original
map $a$ extends $a$ to $U_n$. The base case is a point. The converse
follows from $\pi_1(\operatorname{Spec}R)=1$, and uniqueness follows
from separatedness.

Apply this criterion to an explicit arc at the chosen geometric point
$\bar x$ over $x$. Let $r=|I(x)|$ and $d=\dim U$. Choose branch
parameters $z_1,\ldots,z_r$ in the strict henselian local ring and
complete them to a regular system of parameters. Normal crossings and
smoothness identify its completion with
\[
 \widehat{\mathcal O}_{X_{\overline K},\bar x}
   \simeq\overline K[[z_1,\ldots,z_d]],
 \qquad X\setminus U:\ z_1\cdots z_r=0.
\]
The continuous local homomorphism to $R=\overline K[[t]]$ defined by
\[
 z_i\longmapsto t\quad(1\le i\le r),\qquad
 z_i\longmapsto0\quad(r<i\le d)
\]
defines an arc $a:\operatorname{Spec}R\to X_{\overline K}$ through
$\bar x$. Its generic point lands in $U_{\overline K}$ because the
boundary equation pulls back to the nonzero element $t^r$.
Since $R$ is strictly henselian, the arc factors through the strict
henselization at $\bar x$. On punctured neighbourhoods the covers
$z_i^{1/n}$ pull back to $t^{1/n}$ for every $i$ and every $n$.
Consequently the induced geometric homomorphism is precisely
\[
 \pi_1(\operatorname{Spec}\overline K((t)))
   =T\xrightarrow{d_x}T^{I(x)}\xrightarrow{j_x}\pi_1(U_{\overline K}).
\]
If $j_xd_x$ were zero, the extension criterion just proved would give
$\widetilde a:\operatorname{Spec}R\to U_{\overline K}$ with the same
generic-point map as $a$. Its composite with $U_{\overline K}\to
X_{\overline K}$ equals $a$: two maps from this integral trait to the
separated scheme $X_{\overline K}$ which agree at the generic point
agree everywhere. Their closed point would then lie both in
$U_{\overline K}$ and at $\bar x\in X_{\overline K}\setminus
U_{\overline K}$, a contradiction. Thus $\operatorname{im}(j_xd_x)$
is nonzero. The previously proved equivariance now gives the asserted
nontrivial subgroup of the cyclotomic centralizer.

For the failure of injectivity, take
$U=(\mathbf P^1\setminus\{0,1,\infty\})^2$, compactify in
$(\mathbf P^1)^2$, and blow up $(0,2)$. With old local coordinates
$t$ and $u=z_2-2$, the old boundary is $t=0$. In the chart
$t=ab$, $u=b$, the new boundary has the two branches $a=0$ and $b=0$.
Each branch loop maps to the loop around $t=0$, while the loop around
$u=0$ contracts in the original neighbourhood. Thus the actual map is
\[
 T^2\longrightarrow I_0\subset\pi_1(U_{1,\overline K}),
 \qquad(m,n)\longmapsto m+n.
\]
Its kernel is $\{(m,-m):m\in T\}$ and its image is $I_0\ne1$.
The diagonal maps by multiplication by $2$, so it is indeed nonzero.
The finite matrix reductions in \texttt{code/root\_math\_checks.py}
record this kernel and image modulo $2,3,5,25$.
\end{proof}

The converse requires a geometric boundary centre and a factorization
through its actual punctured henselian group for every section with
nonzero cyclotomic centralizer, with a specified compactification and
arithmetic domain. Merely recognizing a nontrivial abelian subgroup
does not construct that factorization. This obligation already contains
the curve section-surjectivity problem through explicit product maps.
Indeed, let $C/K$ be any hyperbolic curve, let $Q=\mathbf
P^1\setminus\{0,1,\infty\}$, and fix a rational cusp section $t$ of
$Q$. For every section $\sigma$ of $C$, the product section
$(\sigma,t)$ of $C\times Q$ has a centralizer containing
$\{1\}\times I_0$. A boundary-local factorization of that section in
the product compactification would, by the first projection, express
$\sigma$ either as the point section at an interior first coordinate
or as a cuspidal section at a boundary first coordinate. Thus a
universal boundary-factorization assertion for these products entails
the point-or-cusp assertion for all such $\sigma$ over the same field.
For a statement about germs, the same argument retains the corresponding
finite-extension domain. No such exhaustion of the section set is
asserted here.

For source comparison, Carlson--Corwin
\cite[Theorem 3.7 and Proposition 4.12]{B_CarlsonCorwin2026} prove
compactification independence of the \emph{closure} of the boundary
section set and an inclusion into the \emph{closure} of the inductively
defined cuspidal set of a tower. The closure bars are essential in both
statements. Their Proposition 3.4 supplies the precise local group and
action used above; the nonvanishing argument here additionally uses
Hoshi's valuative lemma. None of these comparisons supplies the
universal factorization in the preceding paragraph.

\subsection{Admissible germs and loop recognition}

\paragraph{Statement C9 (Proved; the precise object in scan 66).}\label{res:C9}
The admissible germs for $U_{03}=\mathbf P^1\setminus\{0,1,\infty\}$, defined by trivial cyclotomic centralizer, are exactly the noncuspidal germs. There is a natural injective $\Gamma_\Q$-equivariant map from $U_{03}(\overline\Q)$ to this germ set. Equality is precisely the section-surjectivity assertion for $U_{03}$ over all number fields. Classes of sections defined over all of $\Gamma_\Q$, by themselves, do not give the germ set appearing in this equivariant formulation.
\passage{Section 3, manuscript p.~32, scan 66.}
\begin{proof}
Take the filtered colimit, under restriction, over finite number fields in $\overline\Q$ of the sets of section classes. A point lies over a finite field, so it determines an element of this colimit. If two points have equal germs, their classes agree over some common finite field; the point injectivity in C2 proves equality of the points. Conjugating fields, points and sections by $\Gamma_\Q$ makes this construction equivariant. The identification of admissible and noncuspidal germs is \Rtwo\ applied after restriction. The exhaustion statement follows exactly as in C7, with the quantifier over all fields retained.

There is a concrete type check on the literal whole-group reading. For a section $s:\Gamma_\Q\to E_{03}$, transport by $\tau\in\Gamma_\Q$ using the lift $s(\tau)$ changes the section by
$s(\tau)s(\tau^{-1}\sigma\tau)s(\tau)^{-1}=s(\sigma)$.
Thus this transport on whole-group section classes is trivial. The action on $U_{03}(\overline\Q)$ is nontrivial, for example on the two roots of $z^2-z+1$. Germs carry the required action because their arithmetic domains change. The phrase ``à \ldots\ cas près'' is not used as a theorem: neither the suggested reading ``deux'' nor an unspecified exceptional-case reading repairs a mismatch of arithmetic domains. The exceptions $(0,1)$ and $(0,2)$ calculated in C3 concern different types; the scan here fixes type $(0,3)$, so those calculations do not identify its proposed exceptions. The available mathematics does not identify the omitted words.
\end{proof}

\paragraph{Statement C10 (Proved correction; the character-free strengthening remains unresolved).}\label{res:C10}
For a hyperbolic curve over a number field, its lifted inertia subgroups are exactly the nontrivial subgroups $C(s)$ as $s$ ranges over germs. They are also exactly the maximal procyclic subgroups of $\pi_U$ whose arithmetic normalizer has open image and acts on them through the cyclotomic character. The condition involving only open image of the normalizer, with arbitrary closed subgroups permitted, is false. The same condition restricted to maximal procyclic subgroups without specifying the character is not decided here in the full group.
\passage{La Longue Marche, \S34, cote 140-2, scan 483; cote 149, scans 55--57.}
\begin{proof}
Every geometric cusp and a nonzero tangent vector at its image are defined over a finite extension. C1 supplies the corresponding tangential section, and \Rtwo\ gives $C(s)=I_c$. The reverse inclusion is exactly \Rtwo. The second characterization is Stix~\cite[Theorem 20 and Corollary 21]{Stix2008}, with the arithmetic vanishing of Part A verifying its hypotheses; the maximality step is also explicit: a larger procyclic group would be abelian and intersect $I_c$ nontrivially, and Stix's Corollary 13 then puts that larger group inside $I_c$.

For the unqualified normalizer condition, $\pi_U$ itself and every characteristic subgroup of it have normalizer $E_U$. To make this an explicit counterexample, take $U_{03}$. Its geometric group maps onto the nonabelian finite group $S_3$ by sending its two free generators to $(12)$ and $(123)$. Therefore it is not procyclic, whereas it satisfies the full-image normalizer condition and contains every proper candidate. The closure of its commutator subgroup is characteristic as well. This supplies the comparison map witnessing why the original maximal-subgroup test cannot select the cusp groups. Part E verifies the CM class-two counterexample for the procyclic formulation. The maps from the full group onto these nilpotent quotients send normalizers into quotient normalizers; lifting a quotient normalizer element or a normalized cyclic subgroup back through every level is an additional requirement. The precise remaining full-group question is the existence of a maximal procyclic subgroup with open arithmetic normalizer and no nonzero potentially cyclotomic primary component, already for a once-punctured CM elliptic curve. The finite and Lie calculations below do not supply such a compatible lift.
\end{proof}
